\begin{table}[H]\centering
\caption{Eligibility Paths among Women Observed Before and After the Reform}
\label{tab:identification}\small
\begin{tabular}{lcc}
\toprule
Woman's eligibility path (observed pre and post) & Persons & Share \\
\midrule
Stable young-child across reform & 12,642 & 60.7\% \\
Stable comparison across reform & 5,528 & 26.5\% \\
Left eligibility (child aged out) & 1,313 & 6.3\% \\
Other / multiple switches & 1,093 & 5.2\% \\
Entered eligibility (birth) & 265 & 1.3\% \\
\midrule
\multicolumn{3}{@{}l}{\textit{DiD re-estimated on stable-status women only:} -0.24$^{}$ (0.53) pp} \\
\bottomrule\end{tabular}
\par\vspace{3pt}\footnotesize\raggedright \textit{Notes:} Among women observed both before and after the reform (2022Q2), the table classifies each by how her eligibility (a young child in the household) moves over the panel. Most retain a stable eligibility status across the reform, whose pre/post contrast is uncontaminated by a birth or a child ageing out; the table reports person counts, not the econometric weight each path carries in the pooled estimator. The final row re-estimates Eq.~\eqref{eq:did} on the stable-status women only, as a check that the pooled null is not driven solely by the switchers. The person counts and shares are unweighted; the final-row regression is weighted by the survey weights. Standard errors are clustered at the household level in parentheses. Significance levels: *** $p<0.01$, ** $p<0.05$, * $p<0.10$.
\end{table}
